\section[Propiedades de potencias naturales]{Propiedades de las potencias con exponentes naturales}

Algunas de estas propiedades ya fueron presentadas para introducir definiciones como \(a^0\). Sin embargo aquí se presentarán las propiedades generales.

Las propiedades de esta sección no son reglas aisladas. Todas pueden comprenderse escribiendo las potencias como productos y observando qué ocurre con sus factores.

\subsection{Producto de potencias de igual base}

Considere
\[
2^3\cdot2^4.
\]
Al desarrollar ambas potencias obtenemos
\[
(2\cdot2\cdot2)\cdot(2\cdot2\cdot2\cdot2).
\]
En total aparecen siete factores iguales a \(2\). Por lo tanto,
\[
2^3\cdot2^4=2^7=2^{3+4}.
\]
De manera general, si \(m,n\in\mathbb N\),
\[
a^m\cdot a^n=a^{m+n}.
\]
Si alguno de los exponentes es cero, se requiere \(a\ne0\), pues hemos definido \(a^0\) solamente para bases no nulas.

Observe qué cambia y qué se conserva: la base común permanece; los exponentes se suman porque estamos reuniendo dos grupos de factores iguales.

\subsection{Cociente de potencias de igual base}

Consideremos ahora el cociente
\[
\frac{3^5}{3^2}.
\]
Si desarrollamos las potencias como productos, obtenemos
\[
\frac{3^5}{3^2}
=
\frac{3\cdot3\cdot3\cdot3\cdot3}{3\cdot3}.
\]

Para comprender qué ocurre, podemos utilizar la regla de multiplicación de fracciones estudiada anteriormente y escribir el cociente como un producto:
\[
\frac{3\cdot3\cdot3\cdot3\cdot3}{3\cdot3}
=
\frac33\cdot
\frac33\cdot
\frac31\cdot
\frac31\cdot
\frac31.
\]

Los denominadores iguales a \(1\) no modifican el valor de los últimos factores. En cambio, los dos primeros cocientes valen \(1\):
\[
\frac33=1.
\]
Por lo tanto,
\[
\frac{3^5}{3^2}
=
1\cdot1\cdot3\cdot3\cdot3.
\]

Como \(1\) es el elemento neutro de la multiplicación,
\[
1\cdot1\cdot3\cdot3\cdot3
=
3\cdot3\cdot3
=
3^3.
\]
Así,
\[
\frac{3^5}{3^2}
=
3^3
=
3^{5-2}.
\]

Observe qué ha ocurrido. En el numerador había cinco factores iguales a \(3\), mientras que en el denominador había dos. Cada factor del denominador formó un cociente
\[
\frac33=1
\]
con uno de los factores del numerador. Como esos cocientes iguales a \(1\) no modifican el producto, permanecieron solamente tres factores iguales a \(3\).

Esta es la razón por la que aparece la resta de los exponentes.

El mismo razonamiento puede hacerse con cualquier base no nula. Si \(m>n\), podemos representar
\[
\frac{a^m}{a^n}
\]
como un producto que contiene \(n\) cocientes iguales a
\[
\frac aa=1
\]
y, además, \(m-n\) factores iguales a
\[
\frac a1=a.
\]
Por lo tanto, después de reemplazar los primeros por \(1\), permanecen \(m-n\) factores iguales a \(a\):
\[
\frac{a^m}{a^n}=a^{m-n}.
\]

Si los exponentes son iguales, no queda ningún factor de \(a\). En ese caso,
\[
\frac{a^n}{a^n}=1,
\]
y como ya hemos definido
\[
a^0=1
\qquad(a\ne0),
\]
también podemos escribir
\[
\frac{a^n}{a^n}=a^0=a^{n-n}.
\]

En consecuencia, para \(a\ne0\) y \(m,n\in\mathbb N\) con \(m\ge n\),
\[
\boxed{\frac{a^m}{a^n}=a^{m-n}}.
\]

La condición \(a\ne0\) es necesaria. El razonamiento utiliza cocientes de la forma
\[
\frac aa,
\]
que solamente valen \(1\) cuando \(a\ne0\). Además, necesitamos que el denominador de la expresión original sea distinto de cero.

También pedimos por ahora que
\[
m\ge n,
\]
de modo que \(m-n\) sea un exponente natural. Más adelante, cuando introduzcamos los exponentes negativos, podremos estudiar qué ocurre cuando \(m<n\).

\subsection{Potencia de una potencia}

En la expresión
\[
(2^3)^4,
\]
la base de la potencia exterior es la expresión completa \(2^3\). Desarrollarla significa repetirla como factor cuatro veces:
\[
(2^3)^4=2^3\cdot2^3\cdot2^3\cdot2^3.
\]
Cada grupo aporta tres factores iguales a \(2\), y hay cuatro grupos. En total aparecen \(3\cdot4=12\) factores:
\[
(2^3)^4=2^{12}=2^{3\cdot4}.
\]
Así, para exponentes naturales,
\[
(a^m)^n=a^{m\cdot n},
\]
siempre que las potencias que aparecen estén definidas.

\subsection{Potencia de un producto}

Considere
\[
(2\cdot3)^4.
\]
Al desarrollar la potencia,
\[
(2\cdot3)(2\cdot3)(2\cdot3)(2\cdot3).
\]
Por las propiedades asociativa y conmutativa del producto, podemos reunir los factores iguales:
\[
(2\cdot2\cdot2\cdot2)(3\cdot3\cdot3\cdot3)=2^4\cdot3^4.
\]
De manera general,
\[
(a\cdot b)^n=a^n\cdot b^n.
\]

También, si \(b\ne0\),
\[
\left(\frac ab\right)^n=\frac{a^n}{b^n}.
\]
Esta última igualdad se obtiene multiplicando \(n\) veces la fracción \(a/b\) y aplicando la regla del producto de fracciones como se demostró en la sección anterior.

\subsection{Poner a prueba las propiedades}

Hasta ahora hemos encontrado varias propiedades de las potencias observando cómo se comportan sus factores. Pero es razonable preguntarse hasta dónde llegan esas propiedades.

Por ejemplo, sabemos que
\[
2^3\cdot2^4=2^{3+4}.
\]
La explicación fue sencilla: al desarrollar ambas potencias aparecen tres factores iguales a \(2\) y otros cuatro factores iguales a \(2\). Al juntarlos obtenemos siete factores iguales:
\[
(2\cdot2\cdot2)(2\cdot2\cdot2\cdot2)
=
2^7.
\]

Pero entonces surge una pregunta natural:

\[
\text{¿qué ocurre si las bases son distintas?}
\]

Podríamos preguntarnos, por ejemplo, si es posible hacer algo semejante con
\[
2^3\cdot3^4.
\]

Cuando no estamos seguros de que una operación sea válida, no necesitamos memorizar una prohibición. Podemos formular lo que creemos que podría ocurrir y ponerlo a prueba.

Supongamos, por ejemplo, que intentáramos multiplicar las bases y sumar los exponentes:
\[
2^3\cdot3^4
\stackrel{?}{=}
(2\cdot3)^{3+4}.
\]

El símbolo \(\stackrel{?}{=}\) expresa precisamente nuestra duda: queremos averiguar si ambos lados son realmente iguales.

Calculemos cada uno por separado.

Por un lado,
\[
2^3\cdot3^4
=
8\cdot81
=
648.
\]

Por el otro,
\[
(2\cdot3)^{3+4}
=
6^7
=
279936.
\]

Como
\[
648\ne279936,
\]
la igualdad propuesta es falsa.

Esto nos permite comprender algo más importante que una regla: la propiedad
\[
a^m\cdot a^n=a^{m+n}
\]
funciona porque todos los factores tienen la misma base. Si las bases son diferentes, el razonamiento que dio origen a la propiedad ya no puede hacerse de la misma manera.

Podemos investigar otras dudas siguiendo el mismo procedimiento.

Por ejemplo, alguien podría preguntarse si una potencia puede distribuirse sobre una suma:
\[
(a+b)^n
\stackrel{?}{=}
a^n+b^n.
\]

No necesitamos aceptar ni rechazar la idea sin examinarla. Probemos con valores sencillos:
\[
a=2,\qquad b=3,\qquad n=2.
\]

Entonces, el lado izquierdo es
\[
(2+3)^2
=
5^2
=
25,
\]
mientras que el lado derecho es
\[
2^2+3^2
=
4+9
=
13.
\]

Obtuvimos
\[
25\ne13.
\]

Por lo tanto,
\[
(a+b)^n=a^n+b^n
\]
no puede ser una propiedad válida para todos los números.

Este tipo de ejemplo recibe el nombre de \emph{contraejemplo}. Cuando alguien afirma que una igualdad se cumple siempre, encontrar un solo caso en el que falle es suficiente para demostrar que la afirmación general es falsa.

Hay que tener cuidado con la situación inversa. Encontrar uno, diez o incluso cien ejemplos en los que una igualdad funciona no demuestra que funcionará siempre. Los ejemplos pueden ayudarnos a descubrir un patrón o a sospechar que una propiedad es cierta, pero una propiedad general necesita además una razón que explique por qué se cumple en todos los casos.

Esta distinción será útil durante todo el libro:

\begin{itemize}
    \item si sospechamos que una afirmación es falsa, podemos intentar encontrar un contraejemplo;
    \item si sospechamos que es verdadera, podemos probar algunos casos para comprenderla, pero después necesitaremos justificar por qué funciona siempre.
\end{itemize}

Las propiedades de las potencias no son instrucciones arbitrarias sobre qué símbolos pueden moverse. Cada una expresa algo que ocurre con los números y sus productos. Por eso, cuando aparezca una duda, siempre podemos volver a los significados que ya conocemos, desarrollar las operaciones y dejar que los propios números nos indiquen si nuestra idea puede ser correcta.

\begin{note}[title=\textbf{\color{teal}{Una pequeña pausa filosófica}}]

Hay algo curioso en lo que acabamos de hacer. Planteamos una pregunta matemática y, en cierto sentido, dejamos que las propias matemáticas respondieran.

Preguntamos si
\[
(a+b)^n=a^n+b^n
\]
podía ser siempre cierto. Después elegimos unos valores, seguimos las reglas que ya conocíamos y llegamos a una respuesta que no dependía de nuestra opinión: ambos lados dieron números distintos.

Esta es una de las características más fascinantes de las matemáticas. Nosotros formulamos las preguntas, inventamos símbolos, definiciones y formas de representar las ideas; pero, una vez que aceptamos ciertas reglas, no podemos decidir libremente cuáles serán sus consecuencias. Tenemos que descubrirlas.

Por eso existe una antigua discusión filosófica: ¿las matemáticas se inventan o se descubren?

En este libro adoptaremos una postura sencilla: un poco de ambas.

Podemos crear definiciones, elegir axiomas y construir nuevos sistemas matemáticos. Pero una vez establecidos esos puntos de partida, aparecen relaciones y consecuencias que ya no podemos escoger a nuestro gusto. Podemos explorarlas, intentar comprenderlas y descubrir hasta dónde nos llevan.

En ese sentido, aprender matemáticas no consiste solamente en aprender a manipular símbolos. También consiste en aprender a hacer preguntas y en escuchar las respuestas que surgen al seguir cuidadosamente la lógica.

A veces esas respuestas son inmediatas. Otras veces han requerido años, siglos o todavía no las conocemos.

Y quizá esa sea una de las razones por las que las matemáticas resultan tan profundas: comenzamos jugando con unas pocas ideas y, al seguir sus consecuencias, podemos terminar descubriendo estructuras y patrones que cambian nuestra manera de entender el mundo.
\end{note}

\subsection{Potencias y jerarquía de operaciones}

Al incorporar la potenciación necesitamos ampliar la convención de precedencia. Las potencias se consideran agrupadas antes que los productos, y los productos antes que las sumas. Así,
\[
3+2\cdot5^2
\]
se interpreta como
\[
3+\bigl(2\cdot(5^2)\bigr).
\]
Primero se evalúa \(5^2=25\), luego \(2\cdot25=50\) y finalmente \(3+50=53\).

Esta convención explica la diferencia entre
\[
-2^2=-(2^2)=-4
\]
y
\[
(-2)^2=4.
\]
En la primera expresión, el signo menos no pertenece a la base; en la segunda, los paréntesis indican que la base completa es \(-2\).

Como antes, la precedencia es una convención de escritura que evita paréntesis. No afirma que una operación sea más importante que otra.
